% Draft of a restructured Appendix A (to replace the text from \section{Proofs of the complexity bounds}
% to \end{small} in zetaF.tex). The estimates of Section A.2 are proved in the Lean formalization
% (branch lean-formalization, directory lean/); Lean names are given in typewriter font.
\section{Proofs of the complexity bounds}\label{app:cost}

This appendix proves Theorems~\ref{thm:complexity} and~\ref{thm:uniform} (Sections~\ref{app:tools}--\ref{app:uniform}),
Corollaries~\ref{cor:bernoulli} and~\ref{cor:egf} (Section~\ref{app:sequences}), Theorem~\ref{thm:geometric}(2)
(Section~\ref{app:geometric}), and sketches Claim~C (Section~\ref{app:stretched}). The proof is organized in three
layers: a cost model, stated as assumptions (Section~\ref{app:tools}); a list of estimates, which are pure mathematics
and are proved in the accompanying Lean formalization (Section~\ref{app:estimates}); and an account of the algorithm
that combines them (Sections~\ref{app:pieces} and~\ref{app:uniform}). Theorem~\ref{thm:complexity} is the special
case of Theorem~\ref{thm:uniform} in which all data are fixed, so only the latter is proved.

The following proof text is LLM-generated and unedited by the author.

\begin{small}

\subsection{Cost model}\label{app:tools}
Costs are bit operations on a multitape Turing machine. We use the following standard facts as assumptions.
\begin{itemize}
\item[(M)] $\Mul(n)/n$ is nondecreasing and $\Mul(n\log n) = O(\Mul(n)\log n)$. Then $\Mul(An) = O(\Mul(n))$ for fixed
$A$, and $(\log p)\,\Mul(Ap\log p) = O(\Mul(p)\log^2p)$ (\nolinkurl{regular_scale}, \nolinkurl{regular_cost}).
\item[(K)] Elements of a fixed number field $\mathbb K$ are stored as integer vectors over an integral basis with a
common denominator; an element of height $H$ takes $O(\log H)$ bits, and a product of elements of $b$ bits costs
$O(\Mul(b))$ and has $b + O(1)$ bits.
\item[(BS)] \emph{Binary splitting.} The ordered product of $N$ matrices of fixed size over $\mathcal O_{\mathbb K}$
with entries of $b$ bits, computed by a balanced product tree, costs $O(\Mul(Nb)\log N)$
(\nolinkurl{treeCostL_le}, \nolinkurl{matSize_mul_le}, \nolinkurl{treeCost_matrix_le}). A series whose terms satisfy
$X_{n+1} = q_n^{-1}A_nX_n$, $s_n = s_{n-1} + q_n^{-1}B_nX_n$ is summed by the product tree of the matrices
$\bigl[\begin{smallmatrix}A_n & 0\\ B_n & q_n\end{smallmatrix}\bigr]$ \cite{HP98}.
\item[(A)] Approximate arithmetic on $P$-bit fixed-point numbers costs $O(\Mul(P))$ per operation, with the usual
rounding errors.
\item[(E)] The constants $\pi$ and $\gamma$, and $\exp$ and $\log$ at algebraic points with $O(\log p)$ bits, can be
computed to $p$ bits in $O(\Mul(p)\log^2p)$ \cite{BM80,HP98,BZ10}.
\end{itemize}

\emph{Certified series.} By (M), (BS) and (A), a series is summed to absolute error $2^{-p}$ in
$O(\Mul(Nb)\log N + \Mul(p+G))$ bit operations if it comes with a \emph{certificate} $(N, b, G)$: its terms satisfy a
recurrence of fixed order over $\mathbb K$ whose coefficients have at most $b$ bits for $n \le N$, the tail after $N$
terms is at most $2^{-p}$, and all solutions of the recurrence, normalized as in the product tree, are at most $2^G$
for $n \le N$ (so that $G$ guard bits absorb the rounding errors). With $N = O(p)$, $b = O(\log p)$ and $G = O(p)$
this is $O(\Mul(p)\log^2p)$. Every piece below is a fixed number of certified series, or a bit-burst evaluation.

\emph{Bit-burst.} Let $Y(c) \in \CC^D$ satisfy $cY'(c) = (B_0 + cB_1)Y(c)$, with $B_0, B_1$ over $\mathbb K$, whose
only finite singular point is $c = 0$. Its Taylor coefficients at $c_0 \ne 0$ satisfy
$c_0(l+1)Y_{l+1} = (B_0 + c_0B_1 - l)Y_l + B_1Y_{l-1}$ (\nolinkurl{taylor_rec_iff}). The bit-burst algorithm
\cite{CC90,vdH99} steps from $c_0$ to $c$ through rational points $c_i$ with $2^ib_0$-bit expansions, summing the
Taylor series at $c_i$ in $h = c_{i+1} - c_i$ by binary splitting. The step is controlled by the \emph{Taylor-step
lemma} (\nolinkurl{taylor_step_bound}, \nolinkurl{taylor_step_tail}): if
\[
  \theta \ge |h|\,\frac{\|B_0 + c_iB_1\|}{|c_i|}, \qquad \theta \ge |h|\,\|B_1\|, \qquad \theta \ge \frac{|h|}{|c_i|},
\]
then every solution of the recurrence satisfies $\|h^lY_l\| \le (2\theta)^l\|Y_0\|$; for $\theta < \frac12$ the
step has norm at most $1/(1-2\theta)$ and the tail after $L$ terms is at most $(2\theta)^L\|Y_0\|/(1-2\theta)$.
Take $|c - c_0| \le \frac12|c_0|p^{-A}$, where $p^{A-1} \ge \|B_0\| + |c|\,\|B_1\| + 1$, and
$b_0 \ge 2A\log_2p$. Then $\theta_0 \le \frac12p^{-1}$ for the first step, which therefore has $O(p/\log p)$ terms
with $O(\log p)$-bit coefficients, and $\theta_i \le p^A2^{-2^ib_0} \le 2^{-2^{i-1}b_0}$ for the later ones, which have
$O(p/2^ib_0)$ terms with $O(2^ib_0)$-bit coefficients. Since $\sum_i\theta_i \le \frac14$, the steps amplify errors by at most $e$
(\nolinkurl{error_propagation_linear}), so no a priori bound on $Y$ is needed beyond $|Y(c_0)|$; the total cost is
$O(\Mul(p)\log^2p)$ (\nolinkurl{bitBurst_cost_le}). The convergent Taylor series at $c_i$ continues $Y$ along the
segment $[c_i, c_{i+1}]$.

\subsection{Estimates}\label{app:estimates}
Fix the data of Theorem~\ref{thm:uniform}, let $S \ge 2$ bound $|s_0|$ and $f(x) = \sum a_nx^n$ with
$|a_n| \le CR^n$. The following are proved in the formalization; $A$ denotes a constant depending only on the fixed
data.
\begin{itemize}
\item[(G1)] $|1/\Gamma(w)| \le e^{\pi|\operatorname{Im}w|/2}/\Gamma(\operatorname{Re}w)$ and
$|\Gamma(w)| \le \Gamma(\operatorname{Re}w)$ for $\operatorname{Re}w \ge \frac12$ (\nolinkurl{norm_inv_Gamma_le_G1},
\nolinkurl{norm_Gamma_le_Gamma_re}).
\item[(G2)] $|1/\Gamma(w)| \le e^{1+\pi T/2}(2T+2)^{\lceil T\rceil+1}$ for $|w| \le T$ (\nolinkurl{norm_inv_Gamma_le_G2}).
\item[(G3)] $|\Gamma(w)| \le 2^N\Gamma(\operatorname{Re}w + N)/\min(\delta, 1)$ if $\operatorname{Re}w + N \ge 1$ and
$|w + j| \ge \delta$ for $j < N$ (\nolinkurl{norm_Gamma_le_G3}).
\item[(Li)] \emph{Liouville.} $|\alpha|_v \ge H(\alpha)^{-1}$ for $\alpha \in \mathbb K^\times$ and every archimedean
place $v$, and $|\alpha + \beta|_v \ge (2^{[\mathbb K:\QQ]}H(\alpha)H(\beta))^{-1}$ if $\alpha + \beta \ne 0$
(\nolinkurl{inv_mulHeight}$_1$\nolinkurl{_le_abv}, \nolinkurl{small_divisor}). Hence the divisors $s_0 + \sigma + n$, $a_* + n$ and
the $\delta$ of (G3) are at least $p^{-A}$.
\item[(GS)] $\sum_jx^j|1/\Gamma(\sigma + j + 1)| \le e^{AS\log S}\max(1,x)^{-\operatorname{Re}\sigma}e^x
\le e^{AS\log S}(1+x)^Se^x$ for $|\sigma| \le S$, $x \ge 0$; the tail of
$P(a,y) = y^ae^{-y}\sum_iy^i/\Gamma(a+i+1)$ (\nolinkurl{hasSum_gammaPTerm}) after $N \ge 2y + 2S$ terms is at most
$e^{AS\log S}y^{\operatorname{Re}a}(1+y)^S4^y2^{-N}$ for $|a| \le S$ (the factor $4^y$ cannot be omitted); and
$|\Gamma(w,x)| \le 2x^{\operatorname{Re}w-1}e^{-x}$ for $\operatorname{Re}w \le \frac x2 + 1$
(\nolinkurl{sum_pow_mul_norm_inv_Gamma_le_sharp}, \nolinkurl{sum_pow_mul_norm_inv_Gamma_le},
\nolinkurl{norm_gammaPTerm_tail_uniform}, \nolinkurl{norm_upperGamma_le}).
\item[(KB)] \emph{Kernel.} For $|s| \le S$, $|h_s(u)| \le e^{AS\log S}e^{R^2u}$ for $u \ge u_0$ and
$|h_s(u)| \le e^{AS\log S}u^{-S/2-1}$ for $0 < u \le u_0$; hence $|I_k(s)| \le e^{AS\log S}e^{-(k^2-R^2)u_0}$ for
$k > R$ and $|J_m(s)| \le e^{AS\log S}e^{-\pi^2m^2/(2u_0)}$, for every $s$ with $|s| \le S$, also where the pieces
are defined by continuation (\nolinkurl{norm_kernel_le_far}, \nolinkurl{norm_kernel_le_near},
\nolinkurl{norm_kTerm_le_uniform}, \nolinkurl{norm_dualTerm_le_uniform}). For a twist of period $q$, $\pi^2$ becomes
$\pi^2/q^2$.
\item[(T)] \emph{Truncation.} For fixed $s$, by the bounds in Theorem~\ref{thm:split}, the tails beyond the first
neglected $k$-term and dual term are bounded by geometric series, so $p$ bits need
\begin{equation}\label{eq:counts}
  k \le k_{\max} \approx \sqrt{\frac{p\log 2}{u_0}}, \qquad m \le m_{\max} \approx \frac{q}{\pi}\sqrt{p\log 2\cdot u_0}
\end{equation}
(\nolinkurl{Bounds/Truncation.lean}); uniformly in $|s| \le S$, (KB) gives the same with $p$ replaced by
$p + AS\log S$ and $m_{\max}$ multiplied by $\sqrt2$. Either way there are $O(\sqrt p)$ terms on each side, and the
tails are at most twice their first terms (\nolinkurl{tail_le_two_mul}).
\item[(D)] \emph{Dual moments.} The moments $Z_t[g](c) = \int_0^{u_0}u^{-t}e^{-c/u}g(u)\dd u$ satisfy
\eqref{eq:ibp} and $-\partial_cZ_t = Z_{t+1}$, and $D_a = Z_{1-a}[1]$ satisfies \eqref{eq:Drec}
(\nolinkurl{dualMoment_ibp}, \nolinkurl{hasDerivAt_dualMoment}, \nolinkurl{momentD_rec}); the system of
Proposition~\ref{prop:dual} holds with explicit $B_0, B_1$ (\nolinkurl{dual_system}).
\item[(S)] \emph{Starting values.} $D_{a+n} = \alpha_nD_a + \beta_nE$ with $\alpha_n = (-c)^n/(a+1)_n$ for every
complex $a$ with $a + 1 + i \ne 0$ (\nolinkurl{momentD_add_nat_eq}), and $\sum_nx^n/(n!)^2 \le e^{2\sqrt x}$
(\nolinkurl{sum_pow_div_factorial_sq_le}).
\item[(B)] \emph{Bernoulli numbers.} $B_n(b/f) = -nf^{1-n}L(1-n, \mathbf 1_{b+f\ZZ})$, $f^n(n+1)\#\,B_n(b/f) \in \ZZ$,
$|B_n(x)| \le 4e\cdot n!$ on $[0,1]$, and a real number within $1/(2D)$ of a rational number whose denominator divides
$D$ determines that rational number by rounding (\nolinkurl{bernoulli_eq_LFunction}, \nolinkurl{isInt_den_mul_bernoulli_eval},
\nolinkurl{abs_bernoulli_eval_le}, \nolinkurl{bernoulli_eval_eq_round}).
\end{itemize}

\subsection{The pieces}\label{app:pieces}
By Theorem~\ref{thm:split}, after the reduction of Section~\ref{app:uniform}, the value is a combination of the
$k$-terms $I_k(s)$, $k_0 \le k \le k_{\max}$, the dual terms $J_m(s)$, $1 \le m \le m_{\max}$, and the main term
$M(s)$, at $s = s_0 + \varepsilon$ over $\mathbb K[\varepsilon]/(\varepsilon^{j+1})$. Table~\ref{tab:certificates}
lists how each is computed; we comment on the entries. Throughout, $x = k^2u_0$, $c = \pi^2m^2/q^2$ and
$S = |s_0| + 2$; in Theorem~\ref{thm:complexity}, $S = O(1)$.

\begin{table}[ht]
\centering\footnotesize
\caption{The pieces and their certificates $(N, b, G)$ for $p$ bits (Section~\ref{app:tools}; for the dual terms, the stages of the bit-burst); each costs $O(\Mul(p)\log^2p)$. The guard bits $G$ come from (GS), (KB), (S), (G1)--(G3) and (Li).}
\label{tab:certificates}
\begin{tabular}{@{}lllll@{}}
\toprule
piece & number & $N$ & $b$ & $G$ \\
\midrule
$k$-terms $I_k$, double series \eqref{eq:Ik} & $O(\sqrt p)$ & $O(p)$ & $O(\log p)$ & $O(S\log S + \log p)$ \\
$P(a, x)$ and $1/\Gamma(w)$ & $O(\sqrt p)$, $O(1)$ & $O(p)$ & $O(\log p)$ & $O(p)$ \\
dual terms $J_m$, bit-burst for \eqref{eq:fuchsys} & $O(\sqrt p)$ & stages & stages & $O(p)$ \\
starting values $Y(c_0)$, series \eqref{eq:startvals} & $O(\sqrt p)$ & $O(p)$ & $O(\log p)$ & $O(\sqrt{c_0} + S\log S)$ \\
main term $M$, series in $n$ & $1$ & $O(p)$ & $O(\log p)$ & $O(S\log S + \log p)$ \\
\bottomrule
\end{tabular}
\end{table}

\emph{$k$-terms.}\label{app:kterms} Splitting the series of Theorem~\ref{thm:split} by the parity of $n$ and using
$Q(\sigma_0+m, x) = Q(\sigma_0, x) + e^{-x}\sum_{j<m}x^{\sigma_0+j}/\Gamma(\sigma_0+j+1)$ gives
(\nolinkurl{kTerm_double_series}), for
$f(x) = A_0(x^2) + xA_1(x^2)$ and $\kappa = k^2$,
\begin{equation}\label{eq:Ik}
\begin{split}
  I_k(s) &= k^{-s}\Bigl(Q\bigl(\tfrac s2, x\bigr)A_0\bigl(\tfrac1\kappa\bigr) + Q\bigl(\tfrac{s+1}2, x\bigr)\frac{A_1(1/\kappa)}{k}\Bigr) \\
  &\quad + e^{-x}u_0^{s/2}\sum_{m=1}^{\infty}\frac{a_{2m}}{\kappa^m}\sum_{j=0}^{m-1}\frac{x^j}{\Gamma(\frac s2+j+1)}
   + e^{-x}u_0^{(s+1)/2}\sum_{m=1}^{\infty}\frac{a_{2m+1}}{\kappa^m}\sum_{j=0}^{m-1}\frac{x^j}{\Gamma(\frac{s+1}2+j+1)}.
\end{split}
\end{equation}
Factor out $1/\Gamma(\sigma_0 + j_1 + 1)$, $\sigma_0 \in \{\frac s2, \frac{s+1}2\}$, with $j_1$ the least index for
which $\operatorname{Re}(\sigma_0 + j_1 + 1) \ge \frac12$; the $O(S)$ terms with $j < j_1$ are polynomials in $s$
times that factor, and the others follow a recurrence that divides only by units of
$\mathbb K[\varepsilon]/(\varepsilon^{j+1})$. The double sums are then certified series in $m$ over $\mathbb K$: their
terms are at most $|a_n|\kappa^{-n/2}$ times the sum in (GS), so $O(p/\log(k/R))$ terms and $O(S\log S + \log p)$
guard bits suffice. $A_0(1/\kappa)$ and $A_1(1/\kappa)$ are holonomic values inside the disk of convergence, and
$k^{-s}$ costs one logarithm per $k$.

\emph{Incomplete gamma functions.} For $a \in \ZZ_{\le0}$, $Q(a,y) = 0$. Otherwise $P(a,y)$ is the series of (GS),
with the factor $1/\Gamma(a + i_1 + 1)$ taken out as above; it has $O(p + S)$ terms, and by (GS) its tail is at most
$2^{-p}$ after $N = 2y + 2S + p + O(y + S\log S) = O(p)$ terms, since $y = O(p)$. The
$O(1)$ values $1/\Gamma(w)$ with $\operatorname{Re}w \ge \frac12$, $|w| \le 3S$ come from
$\Gamma(w) = \gamma(w, N) + \Gamma(w, N)$ with $N = O(p)$, where $|\Gamma(w,N)| \le 2N^{3S}e^{-N}$ by (GS) and
$|\Gamma(w)| \ge 0.88\,e^{-3\pi S/2}$ by (G1).

\emph{Dual terms.}\label{app:dual} By Lemma~\ref{lem:kernelhol} and Proposition~\ref{prop:dual},
$J_m(s) = Z_{(3-s)/2}[G_0](c) + Z_{1-s/2}[G_1](c)$ is a fixed linear form in the solution $Y(c)$ of a system
\eqref{eq:fuchsys} of fixed dimension. We multiply $G_0$, $G_1$ by $\Gamma(m_2 + \frac s2)$, $\Gamma(m_2 + \frac{s+1}2)$,
with $m_2$ the least integer making the real parts at least $\frac12$, so that their coefficients are rational
functions of $s$, and replace $E$ by $\widetilde E = Eu_0^{1-t}$. The inhomogeneity $b$ of \eqref{eq:Gsys} then has
entries $\Gamma(\sigma_0+m_2)/\Gamma(\sigma_0+m)$, $m < m_1 + r$, times polynomials of bounded degree, with the common
factor $\beta = (\sigma_0 + m_1 + r)_{m_2-m_1-r}$ of degree $O(S)$; we use $\beta C_{t-i}$ and $\beta\widetilde E$ as
components of $Y$ (the equations of this block are homogeneous in it). Then $B_0, B_1$ are polynomials in $s$ of bounded
degree with coefficients in $\mathbb K$, so their entries have $O(\log p)$ bits and norm $p^{O(1)}$, and the bit-burst
of Section~\ref{app:tools} from a rational $c_0 \le c$ with $c - c_0 \le \frac12cp^{-A}$ costs $O(\Mul(p)\log^2p)$;
$c_0$ has $O(\log p)$ bits. The factors $\beta$, $1/\Gamma$ and $u_0^{1-t}$ are computed once.

\emph{Starting values.} By \eqref{eq:startvals} the components of $Y(c_0)$ are series $\sum_jg_jD_{j-t+1}(c_0)$ whose
exponents lie in two classes modulo $\ZZ$. In each class choose a base exponent $a_*$: the least exponent that occurs,
or $a_* = 0$ in the integer class, with $D_0 = E_1(c_0/u_0) = \Ein(c_0/u_0) - \gamma - \log(c_0/u_0)$. By (S),
$D_a = \alpha_aD_{a_*} + \beta_aE$, where going up from $a_*$ divides by $a_* + n$, at least $p^{-A}$ by (Li), and
going down from $0$ divides by $c_0$. Since $|g_j| \le CR^{2j}/|\Gamma(j+\sigma_0)|$ and
$|\alpha_a| \le c_0^n/|(a_*+1)_n|$, (G1)--(G3) and (S) bound the terms of $\sum_jg_j\alpha_{j-t+1}$ by
$e^{AS\log S}e^{2R\sqrt{c_0}}c_0^{AS}$, and $E\beta_a = D_a - \alpha_aD_{a_*}$ gives the same for the $\beta$-series:
$O(\sqrt{c_0} + S\log S) = O(p)$ guard bits suffice. The base values $D_{a_*}(c_0) = c_0^{a_*}(\Gamma(-a_*) - \gamma(-a_*, c_0/u_0))$
need $\Gamma(-a_*) \le 2^{O(p)}$, by (G3), or $E_1$. For the jets, the divisions by $a + \frac\varepsilon2$ are by
units, and $\partial_a^iD_a$ reduce by \eqref{eq:Drec} to jets of the base values.

\emph{Main term.} $M(s)$ is a series in $n$ whose terms are $a_n$ times powers of $u_0^{1/2}$, $1/\Gamma$ at
$\frac{s+n}2$ or $\frac{s+n}2 + 1$, rational functions of $n$ and $s$ with small divisors $s + n - 1$, and, for $k < k_0$,
$k^{-s-n}P(\frac{s+n}2, k^2u_0) = e^{-x}u_0^{(s+n)/2}\sum_ix^i/\Gamma(\frac{s+n}2+i+1)$; in this form $k^{-n}$ has
cancelled and the double series has positive majorants. At a pole $s_0 = 1 - n_1$ one term is singular and is
expanded explicitly.

\emph{Derivatives in $s$.} All pieces are computed over $\mathbb K[\varepsilon]/(\varepsilon^{j+1})$. The recurrences
are polynomial in $s$ except for the factors $1/\Gamma(\sigma_0 + \varepsilon)$, $u_0^{\varepsilon/2}$, $k^{-\varepsilon}$
and $\beta$, computed once (or once per $k$). The jets of $G_0, G_1$ satisfy a block-triangular system of the form
\eqref{eq:Gsys}, and Proposition~\ref{prop:dual} applies with the moments
$Z^{(i)}_t[g] = \int_0^{u_0}u^{-t}(\log u)^ie^{-c/u}g\dd u$, which satisfy
$c\,Z^{(i)}_{t+1}[g] = Eu_0^{1-t}(\log u_0)^ig(u_0) - (1-t)Z^{(i)}_t[g] - Z^{(i)}_t[\vartheta g] - iZ^{(i-1)}_t[g]$
(\nolinkurl{dualMomentLog_ibp}).
This multiplies all costs by a constant depending on $j$.

\subsection{Proof of Theorem~\ref{thm:uniform}}\label{app:uniform}
\begin{proof}[Proof of Theorems~\ref{thm:complexity} and~\ref{thm:uniform}]
Theorem~\ref{thm:complexity} is the case of fixed data: $F \in \RegQ$ is then a fixed $\mathbb K$-linear combination of
a basis as in Theorem~\ref{thm:uniform}, for a number field $\mathbb K$ containing all data, and the hypotheses hold
for $p \ge p_0$. Let $S = |s_0| + 2$, so that $S\log S \le c'p$ when $|s_0| \le cp/\log p$.

\emph{Reduction.} By Lemma~\ref{lem:reduction}, Remark~\ref{rem:twist} and the decomposition of $\chi$ into the
indicators of the $q$ residue classes, the Laurent coefficient is a combination, with coefficients $\lambda_i\chi(a)$
times fixed algebraic numbers, of Laurent coefficients of bounded order of the functions
$\sum_{k\ge k_0,\,k\equiv a}f(1/k)k^{-s}$ at the points $s_0 + \sigma$, and of the terms $F(k)\chi(k)k^{-s}$ for
$k_1 \le k < k_0$. Here $f$ and $\sigma$ range over fixed finite sets, and each $f$ satisfies (H): each
$B_\sigma = x^{-\sigma}\sum_j(\log x)^jf_{\sigma,j}(1/x)$ is obtained from a solution by a polynomial in the
monodromy operator at $t = 1/x = 0$, so the functions $x^{-\sigma}\sum_j(\log x - 2\pi im)^jf_{\sigma,j}(1/x)$,
$0 \le m \le J$, are solutions, and solving for the powers of $m$ exhibits $f_{\sigma,j}$ as a polynomial in
$\log t$ with coefficients $t^{-\sigma}y(1/t)$ of moderate growth. The values $y_i(k)$, $k < k_0$, are computed once
by analytic continuation \cite{CC90,vdH99}. The coefficients have $O(p)$ bits and size $e^{O(p)}$, so each piece is
needed to absolute error $2^{-O(p)}$, and $|k^{-s_0}| \le k_0^S = 2^{O(p)}$. In the odd case of Remark~\ref{rem:twist}
the kernel is that of $h_{s+\delta}$ with the weights $k^\delta$.

\emph{Cost.} By (T) there are $O(\sqrt p)$ pieces, and by Section~\ref{app:pieces} each is a fixed number of certified
series with $N = O(p)$, $b = O(\log p)$, $G = O(p)$, or a bit-burst evaluation, so each costs $O(\Mul(p)\log^2p)$.

\emph{Memory.} The pieces are computed one after another and added to a running sum; the constants ($\pi$, the factors
$1/\Gamma$, $\beta$, the powers of $u_0$, the base values, and $\log k$ as needed) are $O(1)$ numbers of $O(p)$ bits.
A certified series with $Nb = O(p\log p)$ is split into blocks of $\lceil p/b\rceil$ terms; the product tree of a
block, computed depth-first, occupies $O(p)$ bits and costs $O(\Mul(p)\log p)$, and the $O(Nb/p)$ block matrices are
multiplied into a running product rounded to $O(p)$ bits. The normalized partial products are fundamental matrices
of the recurrence, so all its solutions must be at most $2^{O(p)}$: for the Taylor stages this is the Taylor-step
lemma, for the recurrences of $a_n$ it follows from (H), and the $\Gamma$-twists are of first order. Hence $O(p)$ bits
of memory suffice; without blocking the top of the tree has $O(p\log p)$ bits.

\emph{Large $|s_0|$.} If only $|s_0| \le cp$, run the algorithm at precision $P = \lceil p\log_2p\rceil$; for
$p \ge p_0$ the hypotheses hold with $(P, 2c)$, and the bounds become $O(P^{1/2}\Mul(P)\log^2P) =
O(p^{1/2}\Mul(p)\log^{7/2}p)$ and $O(p\log p)$ bits, by (M).

\emph{Geometric factors.} For the last statement, follow Section~\ref{app:geometric} with the modifications above: the
weights $C_k$, $S_k$ have $O(\sqrt p)$ bits and size $2^{O(\sqrt p)}$, the truncation bound gains the factor
$|z|^{-k}$, and for the $O(1)$ dual points with $\operatorname{Re}c_n < 1$ the values come from the series of
Section~\ref{sec:geometric} in the $D_a(c_n)$, bounded by $2^{O(p)}$ with (G1)--(G3) and (Li).
\end{proof}

\emph{Remarks.} The memory bounds assume multiplication in time $\Mul(p)$ and space $O(p)$, as with FFT-based
algorithms. Some bound on $|s_0|$ is needed with real $u_0$: for $|\operatorname{Im}s_0| \gg p$ the cancellation
$e^{\pi|\operatorname{Im}s_0|/4}$ exceeds $2^{O(p)}$ (Remark~\ref{rem:complexu0}), and for $s_0 \to -\infty$ the value
itself may have more than $p$ bits before the point. For rational $s_0$, and for $s_0$ in an imaginary quadratic
field, $|s_0| \le H(s_0)$; in a general number field a bound on $H(s_0)$ does not bound $|s_0|$. The constants depend
on the operator $L$ (Section~\ref{sec:openuniform}).

\subsection{Bernoulli numbers and exponential generating functions}\label{app:sequences}
\begin{proof}[Proof of Corollary~\ref{cor:bernoulli}]
By $B_n(x+1) = B_n(x) + nx^{n-1}$ we may assume $0 < x = b/f' \le 1$. By (B), $B_n(x) = -nf'^{\,1-n}L_1(1-n,
\mathbf 1_{b+f'\ZZ})$, $D_nB_n(x) \in \ZZ$ with $D_n = f'^{\,n}(n+1)\# = e^{O(n)}$, and $|B_n(x)| \le 4e\cdot n!$, so an
approximation of $L_1(1-n, \cdot)$ with absolute error $e^{-An}$ determines $B_n(x)$, and $n\log_2n + O(n)$ bits suffice.
Theorem~\ref{thm:uniform} applies with $F = 1$ ($L = \partial_x$), $\mathbb K = \QQ$, $j = 0$, $s_0 = 1 - n$ and
$p = \Theta(n\log n)$: $H(s_0) \le n \le p$ and $|s_0| = O(p/\log p)$. The cost is $O(p^{1/2}\Mul(p)\log^2p)$, and
$\Mul(n\log n) = O(\Mul(n)\log n)$. For the Euler polynomials,
$E_n(x) = \frac2{n+1}(B_{n+1}(x) - 2^{n+1}B_{n+1}(\frac x2))$; the tangent and Genocchi numbers are
$2^{2n}(2^{2n}-1)|B_{2n}|/(2n)$ and $2(1-2^n)B_n$, and $B_{n,\chi}$ is a fixed combination of the $f^{n-1}B_n(a/f)$.
\end{proof}

\begin{proof}[Proof of Corollary~\ref{cor:egf}]
Since $a_n = n(n-1)\cdots(n-m+1)\,b_{n-m}$ and $b_n = \lambda^nc_n$, with $c_n = n![t^n]R(e^t)$, it suffices to compute
$c_n$. Over a fixed extension of $\mathbb K$ containing the poles, $R(x)$ is a Laurent polynomial plus a combination of
terms $(1 - x/\rho)^{-i}$. A Laurent polynomial contributes $\sum_j\gamma_jj^n$. For $|\rho| > 1$,
$(1 - e^t/\rho)^{-i} = \sum_{k\ge0}\binom{k+i-1}{i-1}\rho^{-k}e^{kt}$ contributes a fixed combination of the
$\Li_{-n-l}(1/\rho)$, $0 \le l < i$; for $|\rho| < 1$, expand $(-\rho e^{-t})^i(1 - \rho e^{-t})^{-i}$ in the same way;
for a root of unity $\rho \ne 1$, $\Li_{-n-l}(1/\rho) = L_1(-n-l,\chi)$ with $\chi(k) = \rho^{-k}$, by continuity on
$|z| \le 1$, $z \ne 1$. Theorem~\ref{thm:uniform} computes each of these $O(1)$ values. Finally
$|a_n| \le C^nn!$, and with $R = P/Q$, $P, Q \in \mathcal O_{\mathbb K}[x, x^{-1}]$, $Q(1) \ne 0$, $d\lambda \in \mathcal O_{\mathbb K}$
and $N = |N_{\mathbb K/\QQ}(Q(1))|$, the expansion $1/Q(e^{\lambda t}) = \sum_i(Q(1) - Q(e^{\lambda t}))^i/Q(1)^{i+1}$ in
the ring of series $\sum e_nt^n/n!$ with $e_n \in d^{-n}\mathcal O_{\mathbb K}$ gives $d^nN^{n+1}a_n \in \ZZ$. Hence
$n\log_2n + O(n)$ bits determine $a_n$.
\end{proof}

\subsection{Geometric factors}\label{app:geometric}
\begin{proof}[Proof of Theorem~\ref{thm:geometric}(2)]
Expanding $\chi$ into additive characters replaces $z$ by the algebraic numbers $ze^{2\pi ib/q}$, so we may take
$\chi = 1$; by Lemma~\ref{lem:reduction} and the derivatives in $s$, we may take $j = 0$ and $f$ satisfying~(H), and
we use Theorem~\ref{thm:geometric}(1). Since $|z|^{-k}e^{-(k^2-R^2)u_0} \le 2^{-p}$ for
$k \ge \sqrt{p\log2/u_0} + O(1)$ and the dual terms are $O(e^{-\operatorname{Re}c_n/u_0})$, $O(\sqrt p)$ terms suffice;
the cancellation between them is at most $e^{(\log|z|)^2/(4u_0)}$. The $k$-terms at $s$ and $s+1$ are computed as in
Section~\ref{app:pieces}, $C_k$ and $S_k$ incrementally, and $L_k(w)$, $H(w)$ are series like the main term.

For the dual terms, Proposition~\ref{prop:dual} holds for complex $c$ by analytic continuation: $J^{(\delta)}_w(c)$ is a
component of the solution of $cY' = (B_0 + cB_1)Y$, continued to $\CC\setminus(-\infty, 0]$. Compute $c_n$ to $p$
bits and choose a Gaussian rational $c_n'$ with $O(\log p)$ bits, $|c_n - c_n'| \le \frac12|c_n|p^{-A}$, and
$\operatorname{Im}c_n' > 0$ if $c_n$ lies on the cut. The base values $D_a(c_n') = (c_n')^a(\Gamma(-a) - \gamma(-a, c_n'/u_0))$
need $\Gamma(-a)$ once and $\gamma(-a, y)$ at the Gaussian rational $y = c_n'/u_0$, an everywhere convergent series with
$O(p)$ P-recursive terms and cancellation $e^{O(|y|)} = 2^{O(p)}$; the series \eqref{eq:startvals} at $c_n'$ is of the
same kind. The bit-burst from $c_n'$ to $c_n$ is controlled by the Taylor-step lemma, which needs only $|Y(c_n')|$
and no bound on a disk; when $c_n$ lies on the cut, the Taylor series at $c_n'$ continue $Y$ along the segment, which
gives the boundary value from above. Each of the $O(\sqrt p)$ terms costs $O(\Mul(p)\log^2p)$.
\end{proof}

\subsection{Stretched exponentials}\label{app:stretched}
We sketch the cost analysis for $\Lambda(s;\tau,r)$ (Proposition~\ref{prop:expsplit}) with $r = r_1/r_2$ in lowest
terms, $0 < r \le 1$, and algebraic $s$ and $\tau > 0$. Write $a_j = (s-rj)/2$; then $a_j$ modulo $\ZZ$ depends only
on $j$ modulo $2r_2$, and within each class the coefficients $c_j = (-\tau)^j/(j!\Gamma(a_j))$ are hypergeometric in
$j$. We use $|\Gamma(a,x)| \le x^{\operatorname{Re}a-1}e^{-x}$ for $\operatorname{Re}a \le 1$, $x > 0$, and
$|1/\Gamma(a)| \le \Gamma(1-\operatorname{Re}a)\cosh(\pi\operatorname{Im}a)/\pi$ for $\operatorname{Re}a \le 0$.

\emph{$k$-terms.} The $j$-th term of $I_k$ is at most
$e^{-k^2u_0}\,O(1)\,(\tau u_0^{-r/2})^j\,\Gamma(1+\frac{rj}2-\frac{\operatorname{Re}s}2)/j!$, which decays like
$(j!)^{r/2-1}$; within each class the terms are P-recursive by $Q(a+1,x) = Q(a,x) + x^ae^{-x}/\Gamma(a+1)$.

\emph{Dual terms.} For $c > 0$ and $\operatorname{Re}a < 0$, $|D_a(c)| \le c^{\operatorname{Re}a}\Gamma(-\operatorname{Re}a)$,
so the $j$-th term of the series for $J_m$ at $c_0$ is at most $O(1)(\tau c_0^{-r/2})^j\Gamma(\frac{rj}2+O(1))^2/j!$.
This decays like $(j!)^{r-1}$ for $r < 1$, and geometrically with ratio $\tau/(2\sqrt{c_0})$ for $r = 1$, which needs
$c_0 > \tau^2$: for $m \le \tau/\pi$ one starts at such a $c_0$ and continues along the real axis. For $1 < r < 2$ the
terms grow like $(j!)^{r-1}$ for every $c_0$; we do not treat this case. With $v = u^{1/(2r_2)}$ the kernel is a power
of $u$ times an entire function $H(v^{-r_1})$ whose coefficients satisfy
$c_{j+2r_2}/c_j = \tau^{2r_2}(a_j-1)(a_j-2)\cdots(a_j-r_1)/(j+1)_{2r_2}$, so $H$ satisfies an equation with
property~(H). The integration by parts of Proposition~\ref{prop:dual}, with $t$ shifted in steps of $1/(2r_2)$, now
raises the moment index, and solving for the highest index gives a linear system in $c$ whose only finite singular
point is $c = 0$, possibly of higher Poincar\'e rank; the bit-burst applies at $c \gg 0$ unchanged.

\emph{Main term.} $M_\tau(1-s) = \Gamma(\frac{1-s}r)\tau^{(s-1)/r}/r$ needs one value of $\Gamma$ at an algebraic point;
the tails $T_w$ are convergent series with P-recursive terms.

Altogether $\Lambda(s;\tau,r)$ costs $p^{3/2+o(1)}$, the extra factors coming only from the $O(\log p)$ variation in the
lengths of the series. For the products of Proposition~\ref{prop:prodsplit} every piece is a double sum of terms
holonomic in $(n, j)$; by closure under definite summation \cite{Zei90} the inner sums are P-recursive in $n$, and
since unwanted solutions grow at most like $(n!)^{O(1)}C^n$, $O(p\log p)$ guard bits suffice, which again gives
$p^{3/2+o(1)}$.

\end{small}
